\chapter{Common state, prolongation tree, and genealogical measure}
\label{chap:continuo-estado-arbol-medida-comun}

\begin{thesisbox}[title={A single enriched state prior to any internal discrimination}]
The object of this chapter is the complete finite state produced by \(\APP\), oriented by TRIT, and transported by \(\TPK\). From it are constructed the groupoid of reversible symmetries, the tree of all admissible prolongations, the multisection of complete histories, and its canonical measure. No specialized domain is introduced: the domain is the total set of states effectively generated by the same \(\APP\)--TRIT--\(\TPK\) mechanism. Nor is any sheet, continuation, or subsequent reading selected.
\end{thesisbox}

\section{Finite data produced by APP}
\label{sec:continuo-comun-app-finito}

For each depth \(k\geq0\), let \(\mathsf A_k\) denote the finite set
of cellular prefixes produced by APP, and let
\begin{equation}
\label{eq:continuo-comun-app-truncamiento}
 \rho^{\APP}_{k+1,k}:\mathsf A_{k+1}\longrightarrow\mathsf A_k
\end{equation}
the erasure of the last refinement.  This erasure conserves the already
constructed prefix.  The neutral cell of APP gives a prolongation
\begin{equation}
\label{eq:continuo-comun-app-degeneracion}
 \eta^{\APP}_k:\mathsf A_k\longrightarrow\mathsf A_{k+1},
 \qquad
 \rho^{\APP}_{k+1,k}\eta^{\APP}_k=I_{\mathsf A_k}.
\end{equation}
We write
\begin{equation}
\label{eq:continuo-comun-emision-neutra-app}
 \varepsilon^0_k(a):a\longrightarrow\eta^{\APP}_k(a)
\end{equation}
for the emission that realizes that prolongation. Thus, its
destination, \(\eta^{\APP}_k(a)\), and the operation that produces it,
\(\varepsilon^0_k(a)\), are kept separate.
\(\eta^{\APP}_k\) is not interpreted as an absence of history: it records a
length-one refinement with neutral emission and materially preserves the
preceding state.

Each \(a\in\mathsf A_k\) simultaneously bears the two sheets of APP,
\begin{equation}
\label{eq:continuo-comun-dos-hojas}
 \operatorname{Sh}_k(a)
 =\bigl(\Sigma_k(a),\Pi_k(a)\bigr)
 \in\mathsf H_k^{\Sigma}\times\mathsf H_k^{\Pi},
\end{equation}
and not a chosen element of the pair.  For an elementary emission
\(\epsilon:a\to a'\), APP further produces two typed Euclidean divisions,
\begin{equation}
\label{eq:continuo-comun-residuo-cociente-app}
 \delta^{\diamond}(\epsilon)
 =\operatorname{res}^{\diamond}(\epsilon)
  +9\operatorname{quo}^{\diamond}(\epsilon),
 \qquad
 \operatorname{res}^{\diamond}(\epsilon)\in\{0,\ldots,8\},
 \quad
 \operatorname{quo}^{\diamond}(\epsilon)\in\mathbb Z,
\end{equation}
where \(\diamond\in\{\Sigma,\Pi\}\).  The label \(\diamond\) merely indicates
which leaf is being read; both divisions remain in the same ledger.

The cellular source and target of \(\epsilon\) determine its oriented
boundary
\begin{equation}
\label{eq:continuo-comun-frontera-elemental}
 \partial\epsilon:=t(\epsilon)-s(\epsilon).
\end{equation}
The neutral emission satisfies
\begin{equation}
\label{eq:continuo-comun-neutral-app-datos}
 \operatorname{res}^{\diamond}(\varepsilon^0_k(a))=0,
 \qquad
 \operatorname{quo}^{\diamond}(\varepsilon^0_k(a))=0,
 \qquad
 \partial\varepsilon^0_k(a)=0.
\end{equation}

\begin{definition}[Directed graph of APP emissions]
\label{def:continuo-comun-grafo-app}
Let \(\mathsf E_k(a)\) be the finite set of APP emissions of length one
that depart from \(a\in\mathsf A_k\).  We conserve the emissions with
multiplicity: two operations with the same target but with different
residue, quotient, or provenance are different elements of
\(\mathsf E_k(a)\).  By \eqref{eq:continuo-comun-app-degeneracion},
\begin{equation}
\label{eq:continuo-comun-out-no-vacio-app}
 \varepsilon^0_k(a)\in\mathsf E_k(a),
 \qquad
 1\leq \#\mathsf E_k(a)<\infty.
\end{equation}
\end{definition}

\section{TRIT Lift: produced genealogical orientation and carry}
\label{sec:continuo-comun-trit-lift}

The signed incidence of an APP emission is an integer
\(\Delta(\epsilon)\). The balanced representative modulo three uniquely
gives
\begin{equation}
\label{eq:continuo-comun-trit-euclid}
 \Delta(\epsilon)
 =\operatorname{or}(\epsilon)+3\kappa(\epsilon),
 \qquad
 \operatorname{or}(\epsilon)\in\{-1,0,+1\},
 \quad
 \kappa(\epsilon)\in\mathbb Z.
\end{equation}
Thus, orientation and carry are not conditions imposed on a history:
they are total functions of the APP data.

To make the composition visible, let
\(\mathsf C_{\epsilon}:\mathsf M_{s(\epsilon)}\to
\mathsf M_{t(\epsilon)}\) be the memory transport induced by the emission.
The native carry of a composition satisfies
\begin{equation}
\label{eq:continuo-comun-cociclo-carry}
 \kappa(\epsilon_2\epsilon_1)
 =\kappa(\epsilon_2)
  +\mathsf C_{\epsilon_2}\kappa(\epsilon_1),
\end{equation}
whenever \(t(\epsilon_1)=s(\epsilon_2)\). For the identity and the neutral emission,
\begin{equation}
\label{eq:continuo-comun-carry-unidad}
 \kappa(I_a)=0,
 \qquad
 \kappa(\varepsilon^0_k(a))=0,
 \qquad
 \mathsf C_{I_a}=I_{\mathsf M_a}.
\end{equation}

\begin{lemma}[Associativity of the carry ledger]
\label{lem:continuo-comun-asociatividad-carry}
For three composable emissions, the evaluation of
\(\kappa(\epsilon_3\epsilon_2\epsilon_1)\) does not depend on the parenthesization.
\end{lemma}

\begin{proof}
Applying \eqref{eq:continuo-comun-cociclo-carry} twice to the left
parenthesization gives
\[
 \kappa(\epsilon_3)
 +\mathsf C_{\epsilon_3}\kappa(\epsilon_2)
 +\mathsf C_{\epsilon_3}\mathsf C_{\epsilon_2}
  \kappa(\epsilon_1).
\]
The right parenthesization produces the same expression because the transport of
memory is functorial:
\(\mathsf C_{\epsilon_3\epsilon_2}
=\mathsf C_{\epsilon_3}\mathsf C_{\epsilon_2}\).
\end{proof}

The oriented reversal of a reversible emission is denoted by
\(\overline\epsilon\).  TRIT acts by
\begin{equation}
\label{eq:continuo-comun-reversion-trit}
 \operatorname{or}(\overline\epsilon)
 =-\operatorname{or}(\epsilon),
 \qquad
 \partial\overline\epsilon=-\partial\epsilon,
 \qquad
 \mathsf C_{\overline\epsilon}=\mathsf C_\epsilon^{-1}.
\end{equation}
The inverse carry is determined, not chosen, by
\begin{equation}
\label{eq:continuo-comun-carry-inverso}
 \kappa(\overline\epsilon)
 =-\mathsf C_\epsilon^{-1}\kappa(\epsilon),
\end{equation}
for \(\kappa(\overline\epsilon\epsilon)=0\).

\section{TPK transport of the \(2\) leaves and of memory}
\label{sec:continuo-comun-tpk-transporte}

TPK chronology is part of the type of each emission and is not
reconstructed a posteriori by counting letters. Here we recover the action of the
ometer of \cref{def:odometro-9-12,prop:proyeccion-fase-longitud}. Let
\begin{equation}
\label{eq:continuo-comun-ciclo-fase-tpk}
 C_9:=\mathbb Z/9\mathbb Z,\qquad
 \widetilde{\mathcal O}_{9}:=\mathbb Z_{\geq0}\times C_9,\qquad
 \mathcal O_{9,12}:=\mathbb Z/12\mathbb Z\times C_9.
\end{equation}
The unrolled phase-and-memory state is
\(\mathsf Q_k=(w_k,r_k)\in\widetilde{\mathcal O}_{9}\); its reduction
\((w,r)\mapsto([w]_{12},r)\) recovers
\(\mathcal O_{9,12}\). The canonical seed of the clock is
\begin{equation}
\label{eq:continuo-comun-semilla-fase-tpk}
 \mathsf Q_0(a_0):=(0,[0]_9)\qquad(a_0\in\mathsf A_0);
\end{equation}
the other choices of origin are their coherent translations by \(C_9\)
and do not change the update law. For
\(\bar r\in\{0,\ldots,8\}\), we set
\begin{equation}
\label{eq:continuo-comun-carry-fase-tpk}
 \chi_9(r):=\left\lfloor\frac{\bar r+1}{9}\right\rfloor.
\end{equation}
Every emission \(\epsilon\) that lleva depth \(k\) to depth \(k+1\)
acts in the odometro by
\begin{equation}
\label{eq:continuo-comun-accion-fase-tpk}
 \sigma_9(w,r):=(w+\chi_9(r),r+[1]_9),\qquad
 \mathsf Q_{k+1}(\epsilon_*\widehat x_k)
 =\sigma_9\mathsf Q_k(\widehat x_k).
\end{equation}
Its phase component is, by typed definition,
\begin{equation}
\label{eq:continuo-comun-fase-emision-tpk}
 \operatorname{ph}(\epsilon):r_k\longmapsto r_k+[1]_9,
 \qquad
 \operatorname{ph}(\epsilon_n\cdots\epsilon_1):
 r_k\longmapsto r_k+[n]_9.
\end{equation}
We write \(g_k:=r_k\in C_9\) for the published phase and
\(q^{(9)}_k:=w_k\) for the vertical memory, and
\(\beta_k:=[w_k]_{12}\) for its duodecaphase reduction. The iteration of
\eqref{eq:continuo-comun-accion-fase-tpk} is
\begin{equation}
\label{eq:continuo-comun-iteracion-fase-tpk}
 \sigma_9^n(w,r)=
 \left(w+\left\lfloor\frac{\bar r+n}{9}\right\rfloor,
       r+[n]_9\right).
\end{equation}
In efecto, the sumandos
\(\chi_9(r),\chi_9(r+[1]_9),\ldots,\chi_9(r+[n-1]_9)\)
cuentan exactly the cruces of the representative \(8\) to the representative
\(0\), that are
\(\lfloor(\bar r+n)/9\rfloor\). In particular,
\begin{equation}
\label{eq:continuo-comun-retorno-fase-avance-memoria}
 \sigma_9^9(w,r)=(w+1,r),\qquad
 g_{k+9}=g_k,\qquad q^{(9)}_{k+9}=q^{(9)}_k+1.
\end{equation}
Even the neutral emission of APP realizes
\(g_k\mapsto g_{k+1}\): it is neutral in residue, quotient, carry, and boundary,
but records that an extension has occurred. The phase word
\begin{equation}
\label{eq:continuo-comun-palabra-fase-tpk}
 \boldsymbol\tau_k^{(N)}:=(g_k,g_{k+1},\ldots,g_N),
 \qquad g_{j+1}=g_j+[1]_9,
\end{equation}
and the memory \(q^{(9)}\) type the phase return and the
state change as distinct objects.

\subsection{Simultaneous Renewal of the Complete Fiber}
\label{sec:continuo-comun-kph}

The preceding odometer conserves the chronology of each extension. The same
nonadic turn also has the internal Markov realization, already constructed
over the APP--TPK catalogue, which simultaneously renews the two coordinates
of the complete fiber. It is essential to conserve both realizations and not to
identify a generic edge of the tree with an emission of the catalogue without
a typing map.

\begin{proposition}[Internal Markov Realization of the Nonadic Connection]
\label{prop:continuo-comun-kph-heredada}
\label{pII-full:thm:seleccion-novena-fase}
Let the APP seed domain, its observable emission, and its multiplicities be
\[
 T_{\min}:\Omega_{\rm seed}\twoheadrightarrow\mathcal U_6,
 \qquad
 m(u):=\#T_{\min}^{-1}(u),
\]
and the factorization produced by the TPK
\[
 \mathcal U_6\simeq\mathcal S_+\times\mathcal S_\times,
 \qquad |\mathcal S_+|=18,\qquad |\mathcal S_\times|=26.
\]
On functions in \(\mathcal U_6\), the two conditional expectations act
by
\[
 (E_+f)(s,e)=\frac1{18}\sum_{s'\in\mathcal S_+}f(s',e),
 \qquad
 (E_\times f)(s,e)=\frac1{26}
 \sum_{e'\in\mathcal S_\times}f(s,e').
\]
For
\[
 \tau=(+1,-1,0,+1,-1,0,+1,-1,0),
 \qquad
 P_{+1}=E_+,\quad P_{-1}=E_\times,\quad P_0=I,
\]
the catalog transfer and its elevation to all seeds are
\begin{align}
 \mathcal K_{\rm ph}\bigl((u,r),(v,r+1)\bigr)
 &=P_{\tau(r)}(u,v),
 \label{eq:continuo-comun-kph-catalogo}\\
 L_0(\omega,\omega')&=\mathbf1_{\{\omega=\omega'\}},\nonumber\\
 L_+(\omega,\omega')
 &=\frac{E_+(T_{\min}\omega,T_{\min}\omega')}
         {m(T_{\min}\omega')},\nonumber\\
 L_\times(\omega,\omega')
 &=\frac{E_\times(T_{\min}\omega,T_{\min}\omega')}
         {m(T_{\min}\omega')},\nonumber\\
 \widehat{\mathcal K}_{\rm ph}
 \bigl((\omega,r),(\omega',r+1)\bigr)
 &=L_{\tau(r)}(\omega,\omega').
 \label{eq:continuo-comun-kph-elevada}
\end{align}
For \(u=T_{\min}\omega\), it holds phase by phase that
\begin{equation}
\label{eq:continuo-comun-kph-lumpability}
 \sum_{\omega':\,T_{\min}\omega'=v}
 \widehat{\mathcal K}_{\rm ph}
 \bigl((\omega,r),(\omega',r+1)\bigr)
 =
 \mathcal K_{\rm ph}\bigl((u,r),(v,r+1)\bigr).
\end{equation}
In particular,
\begin{equation}
\label{eq:continuo-comun-kph-novena-potencia}
 \boxed{\;
 \mathcal K_{\rm ph}^{\,9}\big|_r=E_+E_\times,
 \qquad
 (E_+E_\times)(u,v)=\frac1{468}.
 \;}
\end{equation}
\end{proposition}

\begin{proof}
In a phase \(+\) or \(\times\), summing
\eqref{eq:continuo-comun-kph-elevada} over
\(T_{\min}^{-1}(v)\) exactly cancels the denominator \(m(v)\); in a
zero phase, \(L_0\) projects onto the identity. This proves
\eqref{eq:continuo-comun-kph-lumpability}. The operators \(E_+\),
\(E_\times\), and \(I\) are idempotent and commute, and each appears three
times in the nonadic word. Therefore
\[
 \prod_{r\in C_9}P_{\tau(r)}
 =E_+^3E_\times^3I^3=E_+E_\times.
\]
The composition first renews one coordinate of cardinal \(18\) and then
the other of cardinal \(26\), so that each target receives
\(1/(18\cdot26)=1/468\).
\end{proof}

\begin{remark}[Two internal realizations, without fictitious identification]
\label{rem:continuo-comun-kph-no-identificacion-aristas}
The transfer \(\mathcal K_{\rm ph}\) is the global realization of the
complete catalog, and \(\sigma_9\) is the chronological action on each state
of the tree. One is not substituted for the other. An edge-by-edge equality
would require an additional map
\[
 \operatorname{Cat}_{k,x}:\operatorname{Out}(x)\longrightarrow\mathcal U_6
\]
that verified the corresponding sum identity on fibers. That map
is not used in this volume. The complete monodromy
\eqref{eq:continuo-comun-kph-novena-potencia} proceeds from the
APP--TPK catalogue, while the genealogy of each history remains in the tree and
its ledger.
\end{remark}

TPK assigns to each admissible emission a typed transport of the leaf pair,
\begin{equation}
\label{eq:continuo-comun-transporte-hojas}
 U_\epsilon:
 \mathsf H_{s(\epsilon)}^{\Sigma}\times
 \mathsf H_{s(\epsilon)}^{\Pi}
 \longrightarrow
 \mathsf H_{t(\epsilon)}^{\Sigma}\times
 \mathsf H_{t(\epsilon)}^{\Pi},
\end{equation}
with
\begin{equation}
\label{eq:continuo-comun-transporte-hojas-funtorial}
 U_{I_a}=I,
 \qquad
 U_{\epsilon_2\epsilon_1}=U_{\epsilon_2}U_{\epsilon_1}.
\end{equation}
The codomain of \(U_\epsilon\) remains the complete pair. No
operation in this chapter replaces
\((\Sigma,\Pi)\) by one of its coordinates.

Let \(m\in\mathsf M_{s(\epsilon)}\) be an accumulated memory.  The TPK action
over it is the affine transformation
\begin{equation}
\label{eq:continuo-comun-accion-memoria}
 \mathsf U_\epsilon(m)
 :=\mathsf C_\epsilon m+\kappa(\epsilon).
\end{equation}

\begin{lemma}[TPK Action on Memory]
\label{lem:continuo-comun-accion-memoria}
The transformations \(\mathsf U_\epsilon\) satisfy
\begin{equation}
\label{eq:continuo-comun-accion-memoria-composicion}
 \mathsf U_{I_a}=I,
 \qquad
 \mathsf U_{\epsilon_2\epsilon_1}
 =\mathsf U_{\epsilon_2}\mathsf U_{\epsilon_1}.
\end{equation}
\end{lemma}

\begin{proof}
The identity follows from \eqref{eq:continuo-comun-carry-unidad}.  For the
composition,
\begin{align*}
 \mathsf U_{\epsilon_2}\mathsf U_{\epsilon_1}(m)
 &=\mathsf C_{\epsilon_2}
   \bigl(\mathsf C_{\epsilon_1}m+\kappa(\epsilon_1)\bigr)
   +\kappa(\epsilon_2)\\
 &=\mathsf C_{\epsilon_2\epsilon_1}m
   +\kappa(\epsilon_2)
   +\mathsf C_{\epsilon_2}\kappa(\epsilon_1)\\
 &=\mathsf C_{\epsilon_2\epsilon_1}m
   +\kappa(\epsilon_2\epsilon_1),
\end{align*}
where the last equality is
\eqref{eq:continuo-comun-cociclo-carry}.
\end{proof}

A finite route is a composable word.
\begin{equation}
\label{eq:continuo-comun-ruta}
 \gamma=\epsilon_1\epsilon_2\cdots\epsilon_k.
\end{equation}
Its boundary, transport, carry, and memory are
\begin{align}
 \partial\gamma
 &:=\sum_{j=1}^{k}\partial\epsilon_j,
 \label{eq:continuo-comun-frontera-ruta}\\
 U_\gamma
 &:=U_{\epsilon_k}\cdots U_{\epsilon_1},
 \label{eq:continuo-comun-transporte-ruta}\\
 \kappa(\gamma)
 &:=\kappa(\epsilon_k)
   +\sum_{j=1}^{k-1}
    \mathsf C_{\epsilon_k}\cdots\mathsf C_{\epsilon_{j+1}}
    \kappa(\epsilon_j),
 \label{eq:continuo-comun-carry-ruta}\\
 m(\gamma;m_0)
 &:=\mathsf U_{\epsilon_k}\cdots
       \mathsf U_{\epsilon_1}(m_0).
 \label{eq:continuo-comun-memoria-ruta}
\end{align}
In \eqref{eq:continuo-comun-frontera-ruta}, the interior endpoints
cancel, leaving
\begin{equation}
\label{eq:continuo-comun-frontera-telescopica}
 \partial\gamma=t(\epsilon_k)-s(\epsilon_1).
\end{equation}

\section{The common state and its total domain}
\label{sec:continuo-comun-estado-total}

\begin{definition}[Elementary ledger]
\label{def:continuo-comun-ledger-elemental}
For an emission that appears between two consecutive enriched prefixes,
we denote by
\(\mathsf Q^-(\epsilon)=(w^-(\epsilon),r^-(\epsilon))\) the odometer of the
source prefix and by
\begin{equation}
\label{eq:continuo-comun-ledger-fase-extremos}
 \mathsf Q^+(\epsilon):=\sigma_9\mathsf Q^-(\epsilon)
 =(w^-(\epsilon)+\chi_9(r^-(\epsilon)),r^-(\epsilon)+[1]_9)
\end{equation}
that of the destination prefix. Thus the odometer is not applied to a bare APP
cell. The emission record is
\begin{equation}
\label{eq:continuo-comun-ledger-elemental}
 \lambda(\epsilon):=\bigl(
 s(\epsilon),t(\epsilon),
 \operatorname{res}^{\Sigma}(\epsilon),
 \operatorname{quo}^{\Sigma}(\epsilon),
 \operatorname{res}^{\Pi}(\epsilon),
 \operatorname{quo}^{\Pi}(\epsilon),
 \operatorname{or}(\epsilon),
 \kappa(\epsilon),
 \partial\epsilon,
 \mathsf Q^-(\epsilon),\mathsf Q^+(\epsilon)\bigr).
\end{equation}
The ledger of \(\gamma\) is the word ordered
\begin{equation}
\label{eq:continuo-comun-ledger-ruta}
 \operatorname{Led}(\gamma)
 :=\bigl(\lambda(\epsilon_1),\ldots,
          \lambda(\epsilon_k)\bigr).
\end{equation}
Preserving this word, and not only its sum, distinguishes two histories with the same terminal output.
\end{definition}

\subsection{Complete reader of the TPK state}
\label{sec:continuo-comun-lector-estado-tpk}

The depth \(K\) of TPK blocks and the depth \(k\) of the
extension tree are distinct types. In this subsection, \(K\) denotes
exclusively the depth of the TPK state already generated by the operator
mechanism; it is not identified with \(k\) or with a generic edge.

\begin{definition}[Compact nonadic memory reader]
\label{def:continuo-comun-lector-compacto-nonadico}
The TPK enriched state of depth \(K\) is
\begin{equation}
\label{eq:continuo-comun-estado-tpk-completo}
 v_K=
 \bigl(
 B_K,\boldsymbol\delta_K,\boldsymbol x_K,
 \boldsymbol{\mathcal C}_K,\Sigma_K,
 \omega_K,\ell_K,g(K),\Lambda_K
 \bigr),
 \qquad
 B_K=(u_K^\pi,u_K^e,u_K^\varphi).
\end{equation}
Its compact reader does not add coordinates: it publishes them from that same state,
\begin{equation}
\label{eq:continuo-comun-lector-compacto-nonadico}
 \mathfrak r_9(v_K)
 :=\bigl(B_K,C_K,p_K,c_K,\Sigma_K,W_K,g(K)\bigr),
\end{equation}
where
\begin{equation}
\label{eq:continuo-comun-lector-compacto-componentes}
\begin{aligned}
 C_K&:=\boldsymbol{\mathcal C}_K
      =\bigl(C_K^\chi\bigr)_{\chi\in\{\pi,e,\varphi\}},\\
 C_K^\chi&:=(N_K^\chi,L_K^\chi,r_K^\chi,
             D_K^\chi,A_K^\chi,B_K^\chi),\\
 p_K&:=(B_0,\ldots,B_K),\\
 c_K&:=(\boldsymbol\delta_K,\boldsymbol x_K,
       \kappa(\gamma^{\rm TPK}_{0:K})),
\end{aligned}
\end{equation}
and the oriented shadow of Witt is preserved as
\begin{equation}
\label{eq:continuo-comun-sombra-witt}
 W_K:=
 \bigl((u_K^\chi,u_K^\chi A_W)\bigr)_
       {\chi\in\{\pi,e,\varphi\}},
 \qquad
 A_W=
 \begin{pmatrix}
 0&1&1&1&1&1\\
 1&0&1&2&2&1\\
 1&1&0&1&2&2\\
 1&2&1&0&1&2\\
 1&2&2&1&0&1\\
 1&1&2&2&1&0
 \end{pmatrix},
 \qquad A_W^2=-I_6.
\end{equation}
The route \(\gamma^{\rm TPK}_{0:K}\) is the word of extensions encoded
by \(\Lambda_K\); its prefix of length \(j\) is the one that produces
\(v_j\). The signature is the output of the explicit internal primitive
\begin{equation}
\label{eq:continuo-comun-primitiva-firma-tpk}
 \Sigma_B:
 \mathsf B_{\rm TPK}\times\mathsf C_{\rm TPK}
 \times\mathbb Z[\mathsf A]\times\Omega_{12}\times\mathsf L_{\rm TPK}
 \longrightarrow\mathsf S_{\rm TPK},
 \quad
 (B,c,\partial\gamma,\omega,\ell)\longmapsto
 \Sigma_B(B,c,\partial\gamma,\omega,\ell),
\end{equation}
with \(\Sigma_B(g\cdot-)=g\cdot\Sigma_B(-)\) for each TPK symmetry. Its
explicit starting structure is
\((B_K,c_K,\partial\gamma^{\rm TPK}_{0:K},\omega_K,\ell_K)\). It is not replaced by an
external signature nor inferred retrospectively from a numerical reader.
\end{definition}

The action of step \(K\mapsto K+1\) is written without omitting the deep state.
For each \(\chi\in\{\pi,e,\varphi\}\),
\begin{equation}
\label{eq:continuo-comun-update-hensel-tpk}
 y_K^\chi=x_K^\chi A_H,\qquad
 u_{K+1}^\chi=y_K^\chi\bmod3,\qquad
 x_{K+1}^\chi=\frac{y_K^\chi-u_{K+1}^\chi}{3},
\end{equation}
with
\[
 A_H=
 \begin{pmatrix}
 2&2&2&1&2&1\\
 2&1&2&2&1&1\\
 1&0&0&1&0&2\\
 0&0&1&1&1&0\\
 2&2&2&0&2&1\\
 2&1&2&1&0&0
 \end{pmatrix},
 \qquad \det A_H=1.
\]
If
\[
 \operatorname{ind}_3(u_1,\ldots,u_6)
 :=\sum_{j=1}^{6}\overline u_j\,3^{6-j}
 \in\{0,\ldots,728\},
\]
the update of the prefix and its \(729\) subcylinders is
\begin{align}
 N_{K+1}^\chi
 &=729N_K^\chi+\operatorname{ind}_3(u_{K+1}^\chi),
 &L_{K+1}^\chi&=L_K^\chi+6,
 \label{eq:continuo-comun-update-cilindro-indice}\\
 I_{K,u}^\chi
 &:=
 \left[
  \frac{729N_K^\chi+u}{3^{L_K^\chi+6}},
  \frac{729N_K^\chi+u+1}{3^{L_K^\chi+6}}
 \right),
 &0&\leq u<729.
 \label{eq:continuo-comun-update-cilindro-intervalo}
\end{align}
The decimal cylinder and the associated carry are updated by the two integer
cases of the same operator \(\mathcal C_{3,1000}\). If a new decimal triad is not
published,
\begin{equation}
\label{eq:continuo-comun-update-cilindro-sin-triada}
 r_{K+1}^\chi=r_K^\chi,\qquad D_{K+1}^\chi=D_K^\chi,\qquad
 A_{K+1}^\chi=729A_K^\chi+
 \operatorname{ind}_3(u_{K+1}^\chi)1000^{r_K^\chi},\qquad
 B_{K+1}^\chi=A_{K+1}^\chi+1000^{r_K^\chi}.
\end{equation}
The primitive \(\mathcal C_{3,1000}\) produces the option datum
\[
 d_{K+1,\chi}^{\rm dec}\in
 \{\bot\}\sqcup\{0,\ldots,999\};
\]
in the first case it equals \(d_{K+1,\chi}^{\rm dec}=\bot\). If a triad occurs, we write
\(d_{K+1,\chi}^{\rm dec}=\delta_{K+1}^\chi\in\{0,\ldots,999\}\) and
\begin{equation}
\label{eq:continuo-comun-update-cilindro-con-triada}
 \begin{aligned}
 r_{K+1}^\chi&=r_K^\chi+1,\qquad
 D_{K+1}^\chi=1000D_K^\chi+\delta_{K+1}^\chi,\\
 A_{K+1}^\chi&=729000A_K^\chi
 +\operatorname{ind}_3(u_{K+1}^\chi)1000^{r_K^\chi+1}
 -729\delta_{K+1}^\chi3^{L_K^\chi},\\
 B_{K+1}^\chi&=A_{K+1}^\chi+1000^{r_K^\chi+1}.
 \end{aligned}
\end{equation}
The register \(\boldsymbol\delta_{K+1}\) adds those triads and the two joint
cocycles
\begin{equation}
\label{eq:continuo-comun-update-triadas-carry-conjunto}
 \begin{gathered}
 q_{K+1}=(u_{K+1}^\pi+u_{K+1}^e+u_{K+1}^\varphi)\bmod3,
 \quad
 c_{K+1}^{\rm joint}=\frac{u_{K+1}^\pi+u_{K+1}^e+u_{K+1}^\varphi-q_{K+1}}3,\\
 a_{K+1}=(u_{K+1}^\pi+u_{K+1}^e-u_{K+1}^\varphi)\bmod3,
 \quad
 h_{K+1}^{\rm joint}=\frac{u_{K+1}^\pi+u_{K+1}^e-u_{K+1}^\varphi-a_{K+1}}3,\\
 \boldsymbol\delta_{K+1}
 =\boldsymbol\delta_K\star
 \bigl((d_{K+1,\chi}^{\rm dec})_\chi,
       c_{K+1}^{\rm joint},h_{K+1}^{\rm joint}\bigr).
 \end{gathered}
\end{equation}
Thus \(C_{K+1}\), \(\boldsymbol\delta_{K+1}\), and
\(\boldsymbol x_{K+1}\) have material actions, not merely names.
The other coordinates are simultaneously updated by
\begin{align}
 p_{K+1}&=p_K\star B_{K+1},
 \label{eq:continuo-comun-update-prefijo-tpk}\\
 \kappa(\gamma^{\rm TPK}_{0:K+1})
 &=\kappa(\varepsilon^{\rm TPK}_{K+1})
   +\mathsf C_{\varepsilon^{\rm TPK}_{K+1}}
    \kappa(\gamma^{\rm TPK}_{0:K}),
 \label{eq:continuo-comun-update-carry-lector-tpk}\\
 c_{K+1}^{\rm comp}
 &:=(\boldsymbol\delta_{K+1},\boldsymbol x_{K+1},
       \kappa(\gamma^{\rm TPK}_{0:K+1})),
 \label{eq:continuo-comun-update-carry-compacto-tpk}\\
 \Sigma_{K+1}
 &=\Sigma_B\bigl(
 B_{K+1},c_{K+1}^{\rm comp},
 \partial\gamma^{\rm TPK}_{0:K}
 +\partial\varepsilon^{\rm TPK}_{K+1},
 \omega_{K+1},\ell_{K+1}\bigr),
 \label{eq:continuo-comun-update-firma-tpk}\\
 W_{K+1}
 &=\bigl((u_{K+1}^\chi,u_{K+1}^\chi A_W)\bigr)_\chi,
 &g(K+1)&=g(K)+[1]_9.
 \label{eq:continuo-comun-update-witt-fase-tpk}
\end{align}
Since \(A_H\) is invertible modulo three and \(A_W^2=-I_6\), the residue
and shadow updates do not collapse distinct blocks. Truncation preserves
the complete sequences
\[
 (B_0,\ldots,B_K),\ (C_0,\ldots,C_K),\
 (p_0,\ldots,p_K),\ (c_0,\ldots,c_K),\
 (\Sigma_0,\ldots,\Sigma_K),\ (W_0,\ldots,W_K)
\]
and simultaneously removes its last entry. Hence the reader
\(\mathfrak r_9\) commutes with the erasure of a block.

\begin{proposition}[Universal scope and certified scope of non-rebooting]
\label{prop:continuo-comun-no-reinicio-alcance}
On every horizontal circuit of nine steps, the following are preserved
\[
 g(K+9)=g(K),\qquad
 q_{\rm mem}(K+9)=q_{\rm mem}(K)+1,
\]
and the prefix is prolonged with nine emissions, so the state is not
reset. In the certified joint section
\(\chi\in\{\pi,e,\varphi\}\), \(0\leq K<387\), it also holds that
\begin{equation}
\label{eq:continuo-comun-no-reinicio-certificado}
 (p_K,C_K)\neq(p_{K+9},C_{K+9}),\qquad
 (B_K,\mathbf N_K,W_K)\neq
 (B_{K+9},\mathbf N_{K+9},W_{K+9}),
 \qquad
 \mathbf N_K:=(N_K^\pi,N_K^e,N_K^\varphi).
\end{equation}
\end{proposition}

\begin{proof}
The two universal identities are the exact iteration of
\eqref{eq:continuo-comun-accion-fase-tpk}; the prefix preserves the nine
entries adjoined by \eqref{eq:continuo-comun-update-prefijo-tpk}. For the
joint section, the recurrence
\eqref{eq:continuo-comun-update-cilindro-indice} compares the \(387\) nonadic pairs
of each of the three channels and produces
\eqref{eq:continuo-comun-no-reinicio-certificado}; this is the assertion
certified in
\cref{thm:generacion-monodromica-conjunta-rev7}. Finally,
\(u\mapsto(u,uA_W)\) is injective, so the change of block also preserves
the change of its complete shadow. This last inequality is not extended
to an arbitrary neutral prolongation outside the certified
section.
\end{proof}

\begin{definition}[APP--TRIT--TPK state of depth \(k\)]
\label{def:continuo-comun-estado}
For a route \(\gamma=\epsilon_1\cdots\epsilon_k\) born in an APP
cell \(a_0\), the common enriched state is
\begin{equation}
\label{eq:continuo-comun-estado}
 \widehat x_k(\gamma):=\Bigl(
  a_0,\ldots,a_k;
  \operatorname{Sh}_0,\ldots,\operatorname{Sh}_k;
  \mathbf{res}_k,\mathbf{quo}_k;
  \mathbf{or}_k,\mathbf\kappa_k;
  \mathsf Q_k;g_k;q^{(9)}_k;\beta_k;m_k;\partial\gamma;\gamma;
  \operatorname{Led}(\gamma)
 \Bigr),
\end{equation}
where
\begin{align*}
 \mathbf{res}_k
 &:=\bigl(
  (\operatorname{res}^{\Sigma}(\epsilon_j),
   \operatorname{res}^{\Pi}(\epsilon_j))
  \bigr)_{1\leq j\leq k},\\
 \mathbf{quo}_k
 &:=\bigl(
  (\operatorname{quo}^{\Sigma}(\epsilon_j),
   \operatorname{quo}^{\Pi}(\epsilon_j))
  \bigr)_{1\leq j\leq k},\\
 \mathbf{or}_k
 &:=\bigl(\operatorname{or}(\epsilon_j)\bigr)_{1\leq j\leq k},
 \qquad
 \mathbf\kappa_k
 :=\bigl(\kappa(\epsilon_j)\bigr)_{1\leq j\leq k},\\
 m_k&:=m(\gamma;m_0).
\end{align*}
The intermediate sheets are obtained by
\(\operatorname{Sh}_{j}=U_{\epsilon_j}\operatorname{Sh}_{j-1}\).
\end{definition}

The simultaneous presence of the route and the ledger prevents identifying
distinct words that terminate in the same cell.  The state is not a list of
cardinals: each coordinate has the type and composition equations of the
preceding sections.

\begin{definition}[Total depth finite domain]
\label{def:continuo-comun-dominio-total}
Let
\begin{equation}
\label{eq:continuo-comun-dominio-total}
 \widehat{\mathcal X}^{\rm tot}_k
 :=\left\{
  \widehat x_k(\epsilon_1\cdots\epsilon_k):
  \begin{array}{l}
   a_0\in\mathsf A_0,\\
   \epsilon_j\in\mathsf E_{j-1}(a_{j-1}),\\
   s(\epsilon_j)=a_{j-1},\ t(\epsilon_j)=a_j
  \end{array}
 \right\}.
\end{equation}
No condition from a later publication is added to \eqref{eq:continuo-comun-dominio-total}.
Being an element of the domain means exactly
having been produced by the total actions
\eqref{eq:continuo-comun-residuo-cociente-app},
\eqref{eq:continuo-comun-trit-euclid}, and
\eqref{eq:continuo-comun-accion-fase-tpk}--
\eqref{eq:continuo-comun-accion-memoria}.
\end{definition}

The common truncation deletes the last emission and all and only the entries of the ledger that it produced:
\begin{equation}
\label{eq:continuo-comun-truncamiento-total}
 \rho_{k+1,k}:\widehat{\mathcal X}^{\rm tot}_{k+1}
 \longrightarrow\widehat{\mathcal X}^{\rm tot}_k,
 \qquad
 \rho_{k+1,k}\widehat x_{k+1}
 =\widehat x_k.
\end{equation}
Let \(\varepsilon^0_k(x)\) be the neutral emission
\eqref{eq:continuo-comun-emision-neutra-app} applied to the APP endpoint of
\(x\). The neutral prolongation defines
\begin{equation}
\label{eq:continuo-comun-seccion-total}
 \eta_k:\widehat{\mathcal X}^{\rm tot}_k
 \longrightarrow\widehat{\mathcal X}^{\rm tot}_{k+1},
 \qquad
 \eta_k(x):=(\varepsilon^0_k(x))_*x,
 \qquad
 \rho_{k+1,k}\eta_k=I.
\end{equation}

\begin{proposition}[Finiteness, non-emptiness, and surjectivity]
\label{prop:continuo-comun-total-finito}
For every \(k\), the set
\(\widehat{\mathcal X}^{\rm tot}_k\) is finite and nonempty, and
\(\rho_{k+1,k}\) is surjective.
\end{proposition}

\begin{proof}
APP produces a nonempty finite set at depth zero. At each step,
every state has a finite set of emissions through
\eqref{eq:continuo-comun-out-no-vacio-app}; by induction, there are only finitely many
words of length \(k\). The sequence of neutral prolongations
produces at least one word of each length. Finally,
\eqref{eq:continuo-comun-seccion-total} is a right inverse of
truncation, so the latter is surjective.
\end{proof}

\section{Finite reversible symmetry groupoid}
\label{sec:continuo-comun-groupoide}

The reversible operations of APP, TRIT, and TPK act on the entire state.
Let \(\mathsf U\) be the group of coherent families
\(g=(g_k)_{k\geq0}\) generated by those operations, where
\begin{equation}
\label{eq:continuo-comun-simetria-coherente}
 \rho_{k+1,k}g_{k+1}=g_k\rho_{k+1,k},
 \qquad
 g_{k+1}\eta_k=\eta_kg_k.
\end{equation}
We denote by \(\mathsf U_k\) its image in
\(\operatorname{Sym}(\widehat{\mathcal X}^{\rm tot}_k)\). This image is
finite because it acts on a finite set. The action of \(g_k\) on
\begin{equation*}
 \widehat x=(a_{\leq k};\operatorname{Sh}_{\leq k};
 \mathbf{res},\mathbf{quo};\mathbf{or},\mathbf\kappa;
 \mathsf Q_k;g_k;q_k^{(9)};\beta_k;m;\partial\gamma;\gamma;
 \operatorname{Led}(\gamma))
\end{equation*}
is
\begin{equation}
\label{eq:continuo-comun-accion-reversible}
\begin{aligned}
 g_k\cdot\widehat x:=\bigl(&
 g_k a_{\leq k};
 U_{g_k}\operatorname{Sh}_{\leq k};
 \mathbf{res}(g_k\gamma),\mathbf{quo}(g_k\gamma);
 \mathbf{or}(g_k\gamma),\mathbf\kappa(g_k\gamma);\\
 &\mathsf Q(g_k\gamma);g(g_k\gamma);q^{(9)}(g_k\gamma);
 \beta(g_k\gamma);
 \mathsf U_{g_k}(m);
 (g_k)_*\partial\gamma;
 g_k\gamma;
 \operatorname{Led}(g_k\gamma)\bigr).
\end{aligned}
\end{equation}
All the terms on the right-hand side are recalculated by means of the APP,
TRIT, and TPK rules; they are not copied as independent labels.

\begin{definition}[Common action groupoid]
\label{def:continuo-comun-groupoide}
The finite groupoid of depth \(k\) is
\begin{equation}
\label{eq:continuo-comun-groupoide}
 \mathfrak G_k
 :=\mathsf U_k\ltimes\widehat{\mathcal X}^{\rm tot}_k.
\end{equation}
Its objects are the common states.  An arrow
\((g_k,\widehat x)\) has source \(\widehat x\), target
\(g_k\cdot\widehat x\), identity \((I,\widehat x)\), and inverse
\((g_k^{-1},g_k\cdot\widehat x)\).
\end{definition}

\begin{lemma}[The action preserves all equations of state]
\label{lem:continuo-comun-groupoide-conserva}
The action \eqref{eq:continuo-comun-accion-reversible} preserves the pair of
sheets, the APP divisions, the TRIT orientation, the carry, the phase
odometer, the vertical memory, its dodecaphase reduction, the affine memory, the
boundary, the route, and the ledger, in their respective types.  Moreover,
\begin{equation}
\label{eq:continuo-comun-groupoide-accion}
 I\cdot\widehat x=\widehat x,
 \qquad
 (g_{2,k}g_{1,k})\cdot\widehat x
 =g_{2,k}\cdot(g_{1,k}\cdot\widehat x).
\end{equation}
\end{lemma}

\begin{proof}
The APP divisions are unique, so the action on an emission
redetermines its residues and quotients.  The balanced decomposition
\eqref{eq:continuo-comun-trit-euclid} is likewise unique.  Composition of the
carry is associative by
\cref{lem:continuo-comun-asociatividad-carry}; the action on memory is associative
by \cref{lem:continuo-comun-accion-memoria}; and the boundary is functorial by
\eqref{eq:continuo-comun-frontera-telescopica}.  Finally, the action on
the ordered word induces the action on its ledger entry by entry.
Coherence with the odometer recalculates \(Q\), \(g\), and \(q^{(9)}\) from
the transformed route and preserves
\eqref{eq:continuo-comun-accion-fase-tpk}. These identities prove
\eqref{eq:continuo-comun-groupoide-accion}.
\end{proof}

\section{Finite tree of all prolongations}
\label{sec:continuo-comun-arbol}

For \(x\in\widehat{\mathcal X}^{\rm tot}_k\), let
\(\operatorname{Out}(x)\) be the set of all admissible TPK emissions
of length one from the endpoint of \(x\).  An emission belongs to
\(\operatorname{Out}(x)\) only when its action jointly updates all
coordinates of \eqref{eq:continuo-comun-estado}.  In particular,
\begin{equation}
\label{eq:continuo-comun-out-finito}
 1\leq d(x):=\#\operatorname{Out}(x)<\infty,
 \qquad
 \varepsilon^0_k(x)\in\operatorname{Out}(x).
\end{equation}

\begin{definition}[Truncated prolongation tree]
\label{def:continuo-comun-arbol-finito}
Given \(N\geq k\) and
\(x\in\widehat{\mathcal X}^{\rm tot}_k\), we define
\(\mathscr T_{k,N}(x)\) as the rooted tree whose vertices at height
\(j\), \(0\leq j\leq N-k\), are the words
\begin{equation}
\label{eq:continuo-comun-vertice-arbol}
 h=(\epsilon_{k+1},\ldots,\epsilon_{k+j}),
 \qquad
 \epsilon_{n+1}\in\operatorname{Out}(x_n),
 \qquad
 x_{n+1}=(\epsilon_{n+1})_*x_n,
\end{equation}
with the empty word as root. An edge appends an emission. Two
words reaching the same terminal state remain distinct vertices
if their ledgers differ.
\end{definition}

The last condition is essential: the tree records genealogies and not the
quotient that remembers only the target.  Its terminal level is
\begin{equation}
\label{eq:continuo-comun-terminal-finito}
 \operatorname{Term}_{k,N}(x)
 :=\{h\in\mathscr T_{k,N}(x):|h|=N-k\}.
\end{equation}
Each \(h\) preserves the complete terminal state.
\begin{equation}
\label{eq:continuo-comun-registro-terminal}
 \operatorname{ter}_{N}(h)
 :=\bigl(x_N,m_N,\partial(\gamma h),
          \gamma h,\operatorname{Led}(\gamma h)\bigr).
\end{equation}

\begin{proposition}[Finiteness and non-emptiness at every horizon]
\label{prop:continuo-comun-arbol-finito-no-vacio}
For all \(x\), \(k\), and \(N\geq k\), the tree
\(\mathscr T_{k,N}(x)\) is finite and
\(\operatorname{Term}_{k,N}(x)\neq\varnothing\).
\end{proposition}

\begin{proof}
The root has a finite positive number of successors by
\eqref{eq:continuo-comun-out-finito}.  Applying the same argument at each
level yields, by induction, finiteness at every finite horizon.  The word
formed by neutral emissions
\((\varepsilon^0_k,\varepsilon^0_{k+1},\ldots,
\varepsilon^0_{N-1})\) is a terminal vertex, so
level \eqref{eq:continuo-comun-terminal-finito} is nonempty.
\end{proof}

The erasure of the last emission defines
\begin{equation}
\label{eq:continuo-comun-proyeccion-terminal}
 \pi_{N+1,N}:\operatorname{Term}_{k,N+1}(x)
 \longrightarrow\operatorname{Term}_{k,N}(x).
\end{equation}

\begin{lemma}[Terminal surjectivity]
\label{lem:continuo-comun-terminal-sobreyectivo}
The map \(\pi_{N+1,N}\) is surjective for every \(N\geq k\).
\end{lemma}

\begin{proof}
Let \(h\in\operatorname{Term}_{k,N}(x)\), with terminal state \(v_h\).
Adding the emission \(\varepsilon_N^0(v_h)\) produces
\(h'\in\operatorname{Term}_{k,N+1}(x)\) and
\(\pi_{N+1,N}(h')=h\).
\end{proof}

\section{Nonadic Connection: Phase, Holonomy, and Memory}
\label{sec:continuo-comun-gamma9-total}

A word of nine emissions is not called a holonomy merely because it has
length nine. The horizontal connection determined by
the phase functor is first preserved. For \(r\in C_9\), we define
\begin{equation}
\label{eq:continuo-comun-relacion-horizontal-fase}
 \begin{split}
 \mathscr R_{r,k}:=\bigl\{
 (x,\epsilon,\epsilon_*x):\;&
 x\in\widehat{\mathcal X}^{\rm tot}_k,\quad
 \epsilon\in\operatorname{Out}(x),\\
 &\mathsf Q_k(x)=(w,r),\quad
 \mathsf Q_{k+1}(\epsilon_*x)=\sigma_9(w,r)
 \bigr\}.
 \end{split}
\end{equation}
The last equality does not restrict the domain: it is the total action
\eqref{eq:continuo-comun-accion-fase-tpk} of every elementary TPK emission. The connection simultaneously transports sheet, memory and ledger:
\begin{equation}
\label{eq:continuo-comun-conexion-tpk-generador}
 \nabla_{\rm TPK}(\epsilon):=
 \bigl(\operatorname{ph}(\epsilon),U_\epsilon,
       \mathsf U_\epsilon,\lambda(\epsilon)\bigr).
\end{equation}
Its composition is calculated by means of
\eqref{eq:continuo-comun-transporte-hojas-funtorial},
\eqref{eq:continuo-comun-accion-memoria-composicion}, and the concatenation of the
ledger.

The apex of the nonadic turn is now the set of phase-compatible horizontal lifts.
\begin{equation}
\label{eq:continuo-comun-gamma9-apice}
 \begin{split}
 \mathsf Z^{\Gamma}_{9,k}:=\coprod_{r\in C_9}\bigl\{
 (x,\epsilon_1,\ldots,\epsilon_9):\;&x_0=x,\quad
 x_i=(\epsilon_i)_*x_{i-1},\\
 &(x_{i-1},\epsilon_i,x_i)
 \in\mathscr R_{r+i-1,k+i-1},\quad 1\leq i\leq9
 \bigr\},
 \end{split}
\end{equation}
 where the subscripts of \(r+i-1\) are read in \(C_9\). The summand is
uniquely determined by \(r=g_k(x)\), so the coproduct does not
duplicate witnesses. The source and
target maps are
\begin{equation}
\label{eq:continuo-comun-gamma9-span}
 \widehat{\mathcal X}^{\rm tot}_k
 \xleftarrow{\ s_{9,k}\ }
 \mathsf Z^{\Gamma}_{9,k}
 \xrightarrow{\ t_{9,k}\ }
 \widehat{\mathcal X}^{\rm tot}_{k+9},
 \qquad
 s_{9,k}(x,\boldsymbol\epsilon)=x,\quad
 t_{9,k}(x,\boldsymbol\epsilon)
 =(\epsilon_9\cdots\epsilon_1)_*x.
\end{equation}
The relational holonomy is
\begin{equation}
\label{eq:continuo-comun-gamma9-relacion}
 \Gamma_{9,k}:=
 \{(s_{9,k}z,t_{9,k}z):z\in\mathsf Z^{\Gamma}_{9,k}\}
 =\operatorname{Hol}_{C_9}(\nabla_{\rm TPK}).
\end{equation}

\begin{lemma}[The ninth level of the tree realizes the holonomy]
\label{lem:continuo-comun-gamma9-arbol}
The map that forgets the intermediate states already determined by the emissions induces a bijection
\begin{equation}
\label{eq:continuo-comun-gamma9-apice-arbol}
 \mathsf Z^{\Gamma}_{9,k}
 \xrightarrow{\ \cong\ }
 \coprod_{x\in\widehat{\mathcal X}^{\rm tot}_k}
 \operatorname{Term}_{k,k+9}(x).
\end{equation}
\end{lemma}

\begin{proof}
Each edge of \(\operatorname{Out}(x_i)\) satisfies
\eqref{eq:continuo-comun-accion-fase-tpk}; hence a terminal word of
length nine determines the nine elements of
\eqref{eq:continuo-comun-relacion-horizontal-fase}. Reciprocally, a horizontal lift
of \eqref{eq:continuo-comun-gamma9-apice} is a composable word
of the tree. The two operations conserve the emissions and are
inverse. Thus \eqref{eq:continuo-comun-gamma9-apice-arbol} is a theorem
derived from the connection, not its definition.
\end{proof}

\begin{proposition}[Phase return, memory advance and no restart]
\label{prop:continuo-comun-gamma9-retorno-memoria}
For every \(z=(x,\epsilon_1,\ldots,\epsilon_9)\) of the apex,
\begin{align}
 g(t_{9,k}z)&=g(s_{9,k}z),
 \label{eq:continuo-comun-gamma9-retorno-fase}\\
 q^{(9)}(t_{9,k}z)&=q^{(9)}(s_{9,k}z)+1,
 \label{eq:continuo-comun-gamma9-avance-memoria}\\
 \operatorname{Led}(t_{9,k}z)
 &=\operatorname{Led}(s_{9,k}z)\star
   \bigl(\lambda(\epsilon_1),\ldots,\lambda(\epsilon_9)\bigr).
 \label{eq:continuo-comun-gamma9-ledger}
\end{align}
In particular, the phase identification
\(g(t_{9,k}z)=g(s_{9,k}z)\) does not induce an identification of states: after
transporting both registers to the same type of comparison, their vertical
memories differ by one unit. The published phase returns, not the
enriched state.
\end{proposition}

\begin{proof}
The first two identities are
\eqref{eq:continuo-comun-retorno-fase-avance-memoria}. The third is the
definition of ledger concatenation. Since the vertical memory differs by
one unit, the complete states are distinct even if some additional visible
coordinate could coincide.
\end{proof}

\begin{proposition}[Isometric compression of the same holonomy]
\label{prop:continuo-comun-gamma9-compresion-isometrica}
Let \(U_{\Gamma_9}:\mathcal H_{\rm full}\to\mathcal H_{\rm full}\) be an
isometric realization of the preceding holonomy and let
\(J:\mathcal H_{\rm vis}\to\mathcal H_{\rm full}\) be the isometric inclusion
of the visible register, with \(E:=J^*\) and \(P:=JE\). We define
\begin{equation}
\label{eq:continuo-comun-gamma9-compresion-datos}
 T:=EU_{\Gamma_9}J,\qquad
 \eta:=(I-P)U_{\Gamma_9}J.
\end{equation}
Then
\begin{equation}
\label{eq:continuo-comun-gamma9-compresion}
 \boxed{\;
 U_{\Gamma_9}J=JT+\eta,\qquad
 I-T^*T=\eta^*\eta.
 \;}
\end{equation}
\end{proposition}

\begin{proof}
Inserting \(I=P+(I-P)\) before \(U_{\Gamma_9}J\) gives
\[
 U_{\Gamma_9}J
 =JEU_{\Gamma_9}J+(I-P)U_{\Gamma_9}J
 =JT+\eta.
\]
The two summands belong to orthogonal subspaces. Since
\(U_{\Gamma_9}\) and \(J\) are isometries,
\[
 I=(U_{\Gamma_9}J)^*(U_{\Gamma_9}J)
   =T^*T+\eta^*\eta,
\]
which is the second identity.
\end{proof}

\begin{proposition}[Totality and Equivariance of Holonomy]
\label{prop:continuo-comun-gamma9-total-natural}
The map \(s_{9,k}\) is surjective. The nine neutral prolongations form
the section
\begin{equation}
\label{eq:continuo-comun-gamma9-seccion-neutra}
 \begin{gathered}
 x_0:=x,\qquad
 \varepsilon^0_{k+i}:=\varepsilon^0_{k+i}(x_i),\qquad
 x_{i+1}:=(\varepsilon^0_{k+i})_*x_i\quad(0\leq i<9),\\
 \eta^{(9)}_k(x):=
 (x,\varepsilon^0_k,\ldots,\varepsilon^0_{k+8})
 \in\mathsf Z^{\Gamma}_{9,k},
 \qquad s_{9,k}\eta^{(9)}_k=I.
 \end{gathered}
\end{equation}
Every coherent symmetry \(a\) whose action on the odometer satisfies
\(\mathsf Q\,a=\bar a\,\mathsf Q\) and
\(\bar a\,\sigma_9=\sigma_9\,\bar a\) induces
\[
 a^Z(x,\epsilon_1,\ldots,\epsilon_9)
 :=(a x,a\epsilon_1,\ldots,a\epsilon_9)
\]
and satisfies
\begin{equation}
\label{eq:continuo-comun-gamma9-naturalidad}
 s_{9,k}a^Z=a_ks_{9,k},
 \qquad
 t_{9,k}a^Z=a_{k+9}t_{9,k}.
\end{equation}
\end{proposition}

\begin{proof}
The neutral word exists in every source and satisfies the phase action \eqref{eq:continuo-comun-accion-fase-tpk}; this proves the section and surjectivity. The condition \(\bar a\sigma_9=\sigma_9\bar a\) sends horizontal relations to horizontal relations. Acting entry by entry preserves composition, carry, memory, boundary, and ledger; reading the first or the final state before or after acting gives \eqref{eq:continuo-comun-gamma9-naturalidad}.
\end{proof}

\section{Total survival and complete multisection}
\label{sec:continuo-comun-supervivencia}

\begin{definition}[Produced survival]
\label{def:continuo-comun-supervivencia}
A state \(x\in\widehat{\mathcal X}^{\rm tot}_k\) survives when
\begin{equation}
\label{eq:continuo-comun-supervivencia}
 \operatorname{Term}_{k,N}(x)\neq\varnothing
 \qquad\text{for every }N\geq k.
\end{equation}
We denote by \(\operatorname{Surv}_k\) the set of such states.
\end{definition}

\begin{theorem}[The total domain is survivor]
\label{thm:continuo-comun-total-superviviente}
For every depth,
\begin{equation}
\label{eq:continuo-comun-surv-igual-total}
 \operatorname{Surv}_k=\widehat{\mathcal X}^{\rm tot}_k.
\end{equation}
In particular, survival does not define a subset supposed before constructing the state.
\end{theorem}

\begin{proof}
The inclusion \(\operatorname{Surv}_k\subseteq
\widehat{\mathcal X}^{\rm tot}_k\) belongs to the definition.  The opposite
inclusion is exactly the nonemptiness at every horizon proved in
\cref{prop:continuo-comun-arbol-finito-no-vacio}.
\end{proof}

\begin{definition}[Multisection of complete prolongations]
\label{def:continuo-comun-multiseccion}
The complete multisection over \(x\) is the inverse limit
\begin{equation}
\label{eq:continuo-comun-multiseccion}
 \operatorname{Msect}(x)
 :=\varprojlim_{N\geq k}
 \bigl(\operatorname{Term}_{k,N}(x),\pi_{N+1,N}\bigr).
\end{equation}
Its elements are all the compatible infinite histories that prolong
\(x\).  None of them receives priority within
\eqref{eq:continuo-comun-multiseccion}.
\end{definition}

\begin{theorem}[No material emptiness of the multisection]
\label{thm:continuo-comun-multiseccion-no-vacia}
For every \(x\in\widehat{\mathcal X}^{\rm tot}_k\),
\begin{equation}
\label{eq:continuo-comun-multiseccion-no-vacia}
 \operatorname{Msect}(x)\neq\varnothing.
\end{equation}
Moreover, every finite coordinate of a history in
\(\operatorname{Msect}(x)\) preserves both leaves, orientation, carry,
memory, boundary, path, and ledger of the state that it prolongs.
\end{theorem}

\begin{proof}
The terminal sets are finite and nonempty by \cref{prop:continuo-comun-arbol-finito-no-vacio}, and their bonding maps are surjective by \cref{lem:continuo-comun-terminal-sobreyectivo}. A compatible element can be constructed recursively: one begins at the root and, at every step, retains a prolongation of the preceding prefix. The neutral prolongation guarantees that the process never stops. Compatibility of the coordinates follows from the fact that each edge of the tree acts through the unique state action defined in \eqref{eq:continuo-comun-estado}.
\end{proof}

\section{Canonical measure generated by the tree}
\label{sec:continuo-comun-medida}

The measure is not introduced as an external weight on the outputs.  It is obtained
by counting the admissible emissions at each vertex before choosing a
continuation.

\begin{definition}[Local weight and horizon measure]
\label{def:continuo-comun-medida-finita}
For \(v\in\widehat{\mathcal X}^{\rm tot}_j\) and
\(\epsilon\in\operatorname{Out}(v)\), we set
\begin{equation}
\label{eq:continuo-comun-peso-local}
 p_v(\epsilon):=\frac{1}{d(v)}\in\mathbb Q_{>0}.
\end{equation}
If
\(h=(\epsilon_{k+1},\ldots,\epsilon_N)\in
\operatorname{Term}_{k,N}(x)\) and \(x_j\) is the prefix that precedes
\(\epsilon_{j+1}\), we define
\begin{equation}
\label{eq:continuo-comun-medida-horizonte}
 \mu_{x,N}(h)
 :=\prod_{j=k}^{N-1}p_{x_j}(\epsilon_{j+1})
 =\prod_{j=k}^{N-1}\frac{1}{d(x_j)}.
\end{equation}
\end{definition}

\begin{lemma}[Normalization]
\label{lem:continuo-comun-medida-normalizada}
For every horizon \(N\geq k\),
\begin{equation}
\label{eq:continuo-comun-medida-suma-uno}
 \sum_{h\in\operatorname{Term}_{k,N}(x)}\mu_{x,N}(h)=1.
\end{equation}
\end{lemma}

\begin{proof}
At height one,
\(\sum_{\epsilon\in\operatorname{Out}(x)}d(x)^{-1}=1\).  Assuming the
normalization at height \(N-k\), each terminal prefix \(h\), with final state
\(v_h\), contributes at the next level
\[
 \sum_{\epsilon\in\operatorname{Out}(v_h)}
 \mu_{x,N}(h)p_{v_h}(\epsilon)
 =\mu_{x,N}(h).
\]
Summing over all prefixes proves the identity by induction.
\end{proof}

\begin{proposition}[Projective consistency]
\label{prop:continuo-comun-medida-consistente}
Finite measures satisfy
\begin{equation}
\label{eq:continuo-comun-medida-pushforward}
 (\pi_{N+1,N})_*\mu_{x,N+1}=\mu_{x,N}.
\end{equation}
\end{proposition}

\begin{proof}
For \(h\in\operatorname{Term}_{k,N}(x)\), the mass of its fiber is
\begin{align*}
 \sum_{\pi_{N+1,N}(h')=h}\mu_{x,N+1}(h')
 &=\mu_{x,N}(h)
   \sum_{\epsilon\in\operatorname{Out}(v_h)}p_{v_h}(\epsilon)\\
 &=\mu_{x,N}(h).
\end{align*}
This is precisely the equality of point measures equivalent to
\eqref{eq:continuo-comun-medida-pushforward}.
\end{proof}

The cylinders of the multisection are
\begin{equation}
\label{eq:continuo-comun-cilindro-msect}
 [h]_{x,N}
 :=\{\omega\in\operatorname{Msect}(x):
       \omega|_N=h\}.
\end{equation}

\begin{theorem}[Canonical Measure of Extensions]
\label{thm:continuo-comun-medida-canonica}
There exists a unique measure of cilindros \(\mu_x\) on
\(\operatorname{Msect}(x)\) such that
\begin{equation}
\label{eq:continuo-comun-medida-cilindro}
 \mu_x([h]_{x,N})=\mu_{x,N}(h).
\end{equation}
It is normalized, has support on the whole multisection, and is obtained only from
the finite sets of APP--TRIT--TPK prolongations.
\end{theorem}

\begin{proof}
The consistency of \cref{prop:continuo-comun-medida-consistente} makes \eqref{eq:continuo-comun-medida-cilindro} independent of the horizon at which a cylinder is represented. The normalization follows from \cref{lem:continuo-comun-medida-normalizada}. Two cylinders are either disjoint or one contains the other; hence the finite masses uniquely define an additive measure on the cylinder algebra. The multisection is a limit of finite discrete sets, and the cylinders form its compact basis; the measure extends uniquely to the \(\sigma\)-algebra they generate. Finally, \(p_v(\epsilon)>0\) for every admissible emission, so every nonempty cylinder has positive mass and the support is the entire multisection.
\end{proof}

\section{Naturality under symmetries and truncations}
\label{sec:continuo-comun-naturalidad}

A coherent family \(g=(g_j)_{j\geq0}\in\mathsf U\) transports an
outgoing emission of depth \(j\) by means of
\begin{equation}
\label{eq:continuo-comun-accion-out}
 (g_{j+1},g_j)_*:\operatorname{Out}(x)\longrightarrow
      \operatorname{Out}(g_j\cdot x),
 \qquad
 \epsilon\longmapsto g_{j+1}\epsilon g_j^{-1}.
\end{equation}
The TPK relations make \eqref{eq:continuo-comun-accion-out} a bijection
that preserves multiplicities.  Acting on each entry yields a tree
isomorphism
\begin{equation}
\label{eq:continuo-comun-isomorfismo-arbol}
 g_{*,N}:\mathscr T_{k,N}(x)
 \xrightarrow{\ \cong\ }
 \mathscr T_{k,N}(g_k\cdot x).
\end{equation}

\begin{theorem}[Naturalness of the tree, the multisection, and the measure]
\label{thm:continuo-comun-naturalidad}
For every coherent reversible symmetry \(g=(g_j)_{j\geq0}\),
\begin{align}
 \pi_{N+1,N}g_{*,N+1}
 &=g_{*,N}\pi_{N+1,N},
 \label{eq:continuo-comun-naturalidad-terminal}\\
 g_*\operatorname{Msect}(x)
 &=\operatorname{Msect}(g_k\cdot x),
 \label{eq:continuo-comun-naturalidad-msect}\\
 (g_{*,N})_*\mu_{x,N}
 &=\mu_{g_k\cdot x,N},
 \qquad
 (g_*)_*\mu_x=\mu_{g_k\cdot x}.
 \label{eq:continuo-comun-naturalidad-medida}
\end{align}
\end{theorem}

\begin{proof}
The first equality states that entrywise action commutes with deleting the
final entry. Passing to the inverse limit yields
\eqref{eq:continuo-comun-naturalidad-msect}. Since
\eqref{eq:continuo-comun-accion-out} is bijective,
\(d(g_j\cdot v)=d(v)\); by
\eqref{eq:continuo-comun-medida-horizonte}, a word and its image have the
same weight. This proves the first equality of measures in
\eqref{eq:continuo-comun-naturalidad-medida}. The second follows on the
cylinders from \eqref{eq:continuo-comun-medida-cilindro}, and the cylinders
determine the measure.
\end{proof}

The truncations of the state and those of the tree also form the square.
\begin{equation}
\label{eq:continuo-comun-cuadrado-truncamiento}
\begin{CD}
 \operatorname{Term}_{k,N+1}(x)
 @>{\pi_{N+1,N}}>>
 \operatorname{Term}_{k,N}(x)\\
 @V{\operatorname{ter}_{N+1}}VV
 @VV{\operatorname{ter}_{N}}V\\
 \widehat{\mathcal X}^{\rm tot}_{N+1}
 @>{\rho_{N+1,N}}>>
 \widehat{\mathcal X}^{\rm tot}_{N}.
\end{CD}
\end{equation}
Commutativity is literal: both paths erase the last emission and its
last ledger entry.

\section{Inverse limit \(\widehat\Omega\) of complete histories and common APP–TRIT–TPK enriched state}
\label{sec:continuo-comun-limite}

The surjections \eqref{eq:continuo-comun-truncamiento-total} produce the
object of complete histories
\begin{equation}
\label{eq:continuo-comun-omega}
 \widehat\Omega
 :=\varprojlim_k
 \bigl(\widehat{\mathcal X}^{\rm tot}_k,\rho_{k+1,k}\bigr).
\end{equation}
An element \(\widehat x=(\widehat x_k)_{k\geq0}\) contains, at each
depth, the same pair of transported sheets, the produced orientations and
carries, the accumulated memory, the boundaries, the route, and the complete
ledger.  Its fiber of continuations from any prefix is
\(\operatorname{Msect}(\widehat x_k)\), equipped with
\(\mu_{\widehat x_k}\).

\begin{theorem}[Common enriched state APP--TRIT--TPK]
\label{thm:continuo-comun-propietario}
The system
\begin{equation}
\label{eq:continuo-comun-propietario}
 \mathfrak C^{\rm gen}:=\Biggl(
 \begin{aligned}
 &\{\widehat{\mathcal X}^{\rm tot}_k,\rho_{k+1,k},\eta_k\}_{k\geq0},
   \widehat\Omega,\{\mathfrak G_k\}_{k\geq0},
   \{\mathscr T_{k,N}\}_{N\geq k},\\
 &\operatorname{Msect},\operatorname{Surv},\{\mu_x\},
   \{\mathfrak r_9(v_K)\}_{K\geq0},
   \mathcal K_{\rm ph},\widehat{\mathcal K}_{\rm ph},\\
 &\operatorname{Sh},\operatorname{or},\kappa,
   \mathsf Q,g,q^{(9)},\beta,m,\partial,\gamma,\operatorname{Led},\\
 &\{\mathsf Z^\Gamma_{9,k},s_{9,k},t_{9,k},\Gamma_{9,k}\}_{k\geq0}
 \end{aligned}
 \Biggr)
\end{equation}
is defined for every state produced by APP--TRIT--TPK\@.  It is non-empty,
natural under reversible symmetries, and compatible with refinement.  In
particular:
\begin{enumerate}
 \item the two leaves are transported together;
 \item orientation and carry it calculan by the division TRIT;
 \item the memory obeys the affine TPK action;
 \item the phase obeys the TPK odometer, returns after nine emissions, and the
       vertical memory advances by one unit;
 \item the Markov realization of that turn satisfies
       \(\mathcal K_{\rm ph}^9=E_+E_\times\) and uniformly renews the complete fiber of \(468=18\cdot26\) emissions;
 \item the compact reader conserves block, cylinder, prefix, carry, signature,
       Witt shadow, and phase, with their updates and truncations;
 \item boundary, path, and ledger are composed without erasing the endpoints or the
       order of the emissions;
 \item every state survives and has a nonempty multisection of complete prolongations;
 \item the canonical measure is projective and has support on all those
       prolongations.
\end{enumerate}
\end{theorem}

\begin{proof}
Totality and finiteness follow from \cref{prop:continuo-comun-total-finito}. The sheet, carry, and memory identities are proved in \cref{lem:continuo-comun-asociatividad-carry,lem:continuo-comun-accion-memoria}; the phase, vertical memory, and nonadic holonomy, including its fiber realization and its state reader, are proved in \cref{prop:continuo-comun-gamma9-retorno-memoria,prop:continuo-comun-gamma9-total-natural,prop:continuo-comun-kph-heredada,prop:continuo-comun-gamma9-compresion-isometrica,prop:continuo-comun-no-reinicio-alcance}; the boundary preserves its endpoints by \eqref{eq:continuo-comun-frontera-telescopica}. Total survival and nonemptiness of the multisections are \cref{thm:continuo-comun-total-superviviente,thm:continuo-comun-multiseccion-no-vacia}. The existence, consistency, and support of the measure are \cref{prop:continuo-comun-medida-consistente,thm:continuo-comun-medida-canonica}. Finally, compatibility with the symmetries and truncations is proved in \cref{thm:continuo-comun-naturalidad} and \eqref{eq:continuo-comun-cuadrado-truncamiento}.
\end{proof}

\section{Structural necessity and origin}
\label{sec:continuo-comun-falsadores}

\begin{proposition}[Falsifiers of the common enriched state]
\label{prop:continuo-comun-falsadores}
Any of the following operations replaces
\(\mathfrak C^{\rm gen}\) with another object and blocks the transfer of the
preceding theorems:
\begin{enumerate}
 \item elegir \(\Sigma\) or \(\Pi\) and descartar the another sheet;
 \item conserving the residue and erasing the quotient;
 \item preserving the TRIT sign and erasing its carry;
 \item replacing the ordered ledger with its terminal value;
 \item identifying two words because they reach the same cell;
 \item choosing a continuation and replacing the multisection with it;
 \item declaring survival without exhibiting prolongations to each horizon;
 \item uniformly normalizing only the last level, ignoring the degrees of
       the prefixes and losing projective consistency;
 \item using a subsequent reading to decide which states enter
       \(\widehat{\mathcal X}^{\rm tot}_k\).
\end{enumerate}
\end{proposition}

\begin{proof}
The first six points eliminate coordinates or relations that appear in the
definition of the state, tree, or multisection. The seventh eliminates the proof of
\cref{thm:continuo-comun-total-superviviente}. The eighth may violate
\eqref{eq:continuo-comun-medida-pushforward} when the terminal degrees are not
constant. The last replaces the total domain
\eqref{eq:continuo-comun-dominio-total} with a retrospective selection.
\end{proof}

The APP cells, the pair of sheets, the TRIT lift, the TPK extensions,
the carry cocycle, memory, boundary, route, and preservation of the
register constitute the starting structures. Upon them, the finite groupoid,
the word tree, the multisection, and the normalized local measure are brought
together in a single domain. The theorems of total survival, nonemptiness,
consistency, and naturality certify this joint construction.
